We first verify the displayed block graphon.  The intervals \(I_1,\ldots,I_r\) partition \([0,1]\) up to null endpoints, so
\[
 \int_0^1\psi_k(u)\psi_\ell(u)\,du
 =r|I_k\cap I_\ell|=\mathbf1\{k=\ell\}.
\tag{1}
\]
If \(T_0\) is the integral operator with kernel \(W_0\), then
\[
 (T_0\psi_k)(u)
 =c_0\mathbf1\{u\in I_k\}\int_{I_k}\sqrt r\,dv
 =\frac{c_0}{r}\psi_k(u).
\tag{2}
\]
The range of \(T_0\) is the span of these \(r\) functions.  Thus \(W_0\) has exact rank \(r\), every nonzero eigenvalue equals \(c_0/r\), and the definition of \(c_0\) gives
\[
 c_0/r\geq\underline\lambda,
 \qquad \int_0^1W_0(u,v)\,dv=c_0/r\geq\kappa^{-1}.
\tag{3}
\]
The compatibility inequality \(c_0\leq\kappa\) gives \(0\leq W_0\leq\kappa\), and \({\rm ass:edge-probability}\) gives \(0\leq\rho_nW_0\leq1\).

Writing Bernstein's moment condition for a centered variable \(X\) as
\[
 \mathbb E|X|^q\leq\frac{q!}{2}\mathbb E(X^2)b_\psi^{q-2},
 \qquad q\geq2,
\tag{4}
\]
we have \(\mathbb E\psi_k(U_i)=r^{-1/2}\) and, for \(r\geq2\),
\[
 |\psi_k(U_i)-r^{-1/2}|\leq(r-1)/\sqrt r\leq b_\psi.
\]
Raising this bound to power \(q-2\), multiplying by the square, and taking expectations proves (4), because \(q!/2\geq1\).  For \(r=1\), the centered eigenfunction is zero.  Equations (1)--(4), together with independent uniform latent types, the prescribed graph draw, and independent Bernoulli assignment, verify every graph, rank, regularity, and design condition in \({\rm def:model-class}\).

Set \(s_0=B/4\), and let independent intercepts \(\xi_i\), independent of graph and assignment variables, have density
\[
 p_{s_0}(t)=s_0^{-1}\cos^2\!\left(\frac{\pi t}{2s_0}\right)
 \mathbf1\{|t|\leq s_0\}.
\tag{5}
\]
It integrates to one.  For deterministic \(g\in C^3[0,1]\), define
\[
 f_i^\pm(a,x)=\xi_i\pm g(x).
\tag{6}
\]
If
\[
 \|g\|_\infty\leq3B/4,
 \qquad \max_{1\leq j\leq2}\|g^{(j)}\|_\infty\leq B,
 \qquad \|g^{(3)}\|_\infty\leq L,
\tag{7}
\]
then the response and smoothness envelopes hold, while unit sampling and anonymous interference follow by construction.  Hence both signs in (6) lie in \(\mathcal M_{n,d,L}\), and
\[
 \theta_n^\pm=\pm g'(1/2),
 \qquad \Delta_g=2|g'(1/2)|.
\tag{8}
\]

We establish one experiment-distance bound for all three choices of \(g\).  The square root of (5), extended by zero, is absolutely continuous and satisfies
\[
 \|\{\sqrt{p_{s_0}}\}'\|_2^2=\frac{\pi^2}{4s_0^2}.
\]
The fundamental theorem of calculus and Minkowski's inequality therefore give
\[
 H^2\{p_{s_0}(\mathord\cdot-\eta_1),p_{s_0}(\mathord\cdot-\eta_2)\}
 \leq\frac{\pi^2(\eta_1-\eta_2)^2}{4s_0^2}.
\tag{9}
\]
Let \(X_i^\circ=1/2+V_i\) for \(N_i>0\) and \(X_i^\circ=0\) for \(N_i=0\).  Conditional product affinity, followed by averaging over the common \((A,Z)\) marginal, gives
\[
 H^2(Q_g^+,Q_g^-)
 \leq\frac{\pi^2}{s_0^2}\mathbb E\sum_{i=1}^ng(X_i^\circ)^2,
 \qquad \operatorname{TV}(Q_g^+,Q_g^-)\leq H(Q_g^+,Q_g^-).
\tag{10}
\]

Put \(\bar c=c_0/r\).  A unit connects independently to each other unit with marginal probability \(p_n=\bar c d/n\), hence
\[
 N_i\sim{\rm Binomial}(n-1,p_n),
 \qquad \mu_n=(n-1)p_n\asymp d.
\tag{11}
\]
The binomial identity \(\mathbb E(N_i+1)^{-1}=\{1-(1-p_n)^n\}/(np_n)\), the inequality \(N^{-1}\mathbf1(N>0)\leq2(N+1)^{-1}\), and \(\mathbb E_Z(V_i^2\mid A,N_i=N)=1/(4N)\) imply
\[
 \mathbb E\sum_iV_i^2\mathbf1(N_i>0)=O(n/d),
 \qquad n\Pr(N_i=0)=O(ne^{-cd}).
\tag{12}
\]
Take
\[
 g_G(x)=c_GB\sqrt{d/n}\,(x-1/2).
\tag{13}
\]
Because \(d/n\to0\), (7) holds for all sufficiently large \(n\).  Equations (10)--(13) give \(H^2(Q_G^+,Q_G^-)\leq Cc_G^2\).  With \(\eta=(1-2\alpha)/2>0\), choose a fixed \(c_G>0\) small enough that
\[
 \operatorname{TV}(Q_G^+,Q_G^-)\leq\eta,
 \qquad \Delta_G=2c_GB\sqrt{d/n}.
\tag{14}
\]

For the interior pair, suppose \(L>0\), define
\[
 \phi(t)=t(1-4t^2)^4\mathbf1\{|t|<1/2\},
 \qquad g_h(x)=c_ILh^3\phi\!\left(\frac{x-1/2}{h}\right).
\tag{15}
\]
The zero extension of \(\phi\) is \(C^3\) because of the fourth-order endpoint zeros, and \(\phi'(0)=1\).  If \(C_j=\|\phi^{(j)}\|_\infty\), then \(\|g_h^{(j)}\|_\infty=c_ILh^{3-j}C_j\).  Choosing \(c_IC_3\leq1\), and using \(h\leq d^{-1/2}\to0\), verifies (7).

The largest atom of \({\rm Binomial}(N,1/2)\) is at most \(CN^{-1/2}\), while \(|M/N-1/2|\leq h\) contains at most \(2Nh+2\) lattice points.  Therefore
\[
 \Pr\{|M/N-1/2|\leq h\mid N\}\leq C(\sqrt N\,h+N^{-1/2}).
\tag{16}
\]
For the binomial degree in (11), a Chernoff split at \(\mu_n/2\) and Jensen's inequality give \(\mathbb E\sqrt{N_i}\leq C\sqrt d\) and \(\mathbb E\{N_i^{-1/2}\mathbf1(N_i>0)\}\leq Cd^{-1/2}\).  Since \(h\geq d^{-1}\), (16) implies
\[
 \Pr(|V_i|\leq h,N_i>0)\leq C\sqrt d\,h.
\tag{17}
\]
The bump vanishes for isolates and outside this event, so (10), (15), and (17) yield
\[
 H^2(Q_h^+,Q_h^-)\leq Cc_I^2n\sqrt d\,L^2h^7.
\tag{18}
\]
For \(h_0=\{L^{-2}/(n\sqrt d)\}^{1/7}\) in the interior regime, choose fixed \(c_I\) small enough to obtain
\[
 \operatorname{TV}(Q_{h_0}^+,Q_{h_0}^-)\leq\eta,
 \qquad \Delta_I=2c_ILh_0^2
 =2c_IL^{3/7}(n^2d)^{-1/7}.
\tag{19}
\]

For the lattice pair, fix \(K>e\bar c\) so large that the Chernoff exponent \(a_K\) satisfies \(a_KC_{\log}>2\), and put \(D_n=\lceil Kd\rceil\).  Optimizing exponential Markov's inequality in (11) and taking a union bound gives
\[
 \Pr(\max_iN_i>D_n)\leq ne^{-a_Kd}=o(1).
\tag{20}
\]
Set \(w_n=(4D_n)^{-1}\) and
\[
 g_D(x)=c_DLw_n^3\phi\!\left(\frac{x-1/2}{w_n}\right),
 \qquad c_DC_3\leq1.
\tag{21}
\]
Derivative scaling verifies (7).  If \(1\leq N_i\leq D_n\), every noncentral exposure-lattice point has absolute centered value at least \((2D_n)^{-1}>w_n\); at a central point the odd bump is zero, and at an isolate it is also zero.  Thus the two observed laws coincide on \(\{\max_iN_i\leq D_n\}\), and
\[
 \operatorname{TV}(Q_D^+,Q_D^-)\leq o(1)\leq\eta,
 \qquad \Delta_D=2c_DLw_n^2\asymp Ld^{-2}.
\tag{22}
\]
This proves membership and the asserted linear, smooth-bump, and lattice witnesses inside the explicit block submodel.

For any of these pairs, an estimator \(T\) satisfies the two-point inequality
\[
 \max_{\sigma\in\{+,-\}}\mathbb E_{Q^\sigma}|T-\theta^\sigma|
 \geq\frac\Delta4\{1-\operatorname{TV}(Q^+,Q^-)\}.
\tag{23}
\]
Indeed, split according to \(|T-\theta^+|\geq\Delta/2\) and use the definition of total variation.  If an interval \(\mathcal C\) covers both laws with probability at least \(1-\alpha\), then
\[
 Q^+\{\theta^+,\theta^-\in\mathcal C\}
 \geq1-2\alpha-\operatorname{TV}(Q^+,Q^-),
\]
and hence
\[
 \mathbb E_{Q^+}\operatorname{length}(\mathcal C)
 \geq\Delta\{1-2\alpha-\operatorname{TV}(Q^+,Q^-)\}.
\tag{24}
\]
Equations (14), (19), and (22)--(24) provide the three lower orders for both minimax criteria.

Finally put \(a_n=(n\sqrt d)^{-1/2}\).  The balance relation is
\[
 Lh_0^2=a_nh_0^{-3/2}=L^{3/7}(n^2d)^{-1/7}.
\tag{25}
\]
If \(h_0\leq d^{-1}\), evaluation at \(d^{-1}\) and the equivalent inequality \(L^2\geq d^{13/2}/n\) give
\[
 \inf_{h\in\mathcal H_{n,d}}\{Lh^2+a_nh^{-3/2}\}\leq2Ld^{-2}.
\tag{26}
\]
If \(d^{-1}<h_0<d^{-1/2}\), evaluation at \(h_0\) gives at most twice the quantity in (25).  If \(h_0\geq d^{-1/2}\), then \(L^2\leq d^3/n\), and evaluation at \(d^{-1/2}\) gives at most \(2\sqrt{d/n}\).  For \(L=0\), that same endpoint gives exactly \(\sqrt{d/n}\).  Combining the graph pair with the lattice pair in the first regime, with the interior pair in the second, and alone in the third proves
\[
 R^*_{n,d,L}\geq c\mathfrak r^{\mathrm{sup}}_{n,d,L},
 \qquad \Lambda^*_{n,d,L}(\alpha)\geq c\mathfrak r^{\mathrm{sup}}_{n,d,L}.
\tag{27}
\]
The constant is independent of \(n,d\).  The block construction used the two displayed compatibility inequalities only as sufficient membership conditions.  Independently, \({\rm thm:unrestricted-supported-converse}\) proves (27) over every eventually nonempty repaired class, which establishes the theorem's final sentence without claiming that the block conditions characterize nonemptiness.